%% 第2巻　数と量のことば　本文（第7回編集委員会）
%% 序章 p.8〜11、第1部 p.12〜85、第2部 p.86〜135、第3部 p.136〜209、第4部 p.210〜247、巻末 p.248〜
%% 章は \bkinputchap{ch2-NN} で読む（まだない章は、組むときに飛ばして「BKPAGE missing」を出す）
\providecommand\bkinputchap[1]{\IfFileExists{#1.tex}{\input{#1}}{\clearpage\typeout{BKPAGE missing #1}}}
\bkinputchap{ch2-00}

%% ======================== 第1部 量と数 ========================
\bkPartOpener{1}{\R{第}{だい}2\R{巻}{かん}\hspace{0.8em}\R{第}{だい}1\R{部}{ぶ}}{\R{量}{りょう}と\R{数}{かず}}%
  {172/{「5」と「\qty{5}{m}」は、\R{何}{なに}がちがうのかな？},
   160/{「0\,℃」は、\R{温度}{おん|ど}が「ない」という\R{意味}{い|み}？},
   148/{ピザ2\R{枚}{まい}を3\R{人}{にん}で\R{分}{わ}けると、1\R{人分}{にん|ぶん}は？}}%
  {% ものさし・ストップウォッチ・温度計（白い線）
   \begin{scope}[white, line width=1.2pt, line cap=round, line join=round]
     \draw[rounded corners=1mm] (118,116) rectangle (160,125);
     \foreach \x in {121,124,...,157} \draw[line width=0.7pt] (\x,125) -- (\x,122.5);
     \draw (136,140) circle (8.5); \draw (136,148.5) -- (136,151.5); \draw (133,152) -- (139,152);
     \draw (136,140) -- (140.5,145);
     \draw (151,135) -- (151,160) arc (180:0:2.5) -- (156,135);
     \draw (153.5,130.5) circle (5); \fill (153.5,130.5) circle (2.6);
     \draw[line width=2pt] (153.5,133) -- (153.5,149);
   \end{scope}}%
  {\R{数}{かず}には「かぞえる\R{数}{かず}」と「はかる\R{数}{かず}」があるよ。この\R{部}{ぶ}では、\R{量}{りょう}を\R{数}{かず}で\R{書}{か}くための、いちばん\R{大事}{だい|じ}なきまりを\R{学}{まな}ぶよ。}
\bkPartContents{1}{%% 自動生成（gen.py）
\bktocchap{2-1}{「かぞえる\R{数}{かず}」と「はかる\R{数}{かず}」}{「5」だけでは、\R{何}{なん}のことかわからない}{7}{1}{1}{18}{C}
\bktocchap{2-2}{マイナスの\R{数}{かず}と\R{数直線}{すう|ちょく|せん}}{0\,℃より\R{寒}{さむ}いときは、どう\R{書}{か}く？}{6}{1}{2}{38}{C}
\bktocchap{2-3}{＋ － × ÷ は、\R{何}{なに}を\R{表}{あらわ}しているの？}{かけ\R{算}{ざん}の\R{答}{こた}えは、\R{長方形}{ちょう|ほう|けい}の\R{広}{ひろ}さ？}{1,2}{1}{1}{60}{C}
\bktocchap{2-4}{\R{分数}{ぶん|すう}と\R{小数}{しょう|すう}の\R{計算}{けい|さん}}{ピザ2\R{枚}{まい}を3\R{人}{にん}で\R{分}{わ}けると、1\R{人分}{にん|ぶん}は\R{何枚}{なん|まい}？}{1}{1}{2}{76}{C}
\bktocchap{2-5}{\R{比}{ひ}・\R{割合}{わり|あい}・\R{百分率}{ひゃく|ぶん|りつ}}{「100\,\%」をこえることはある？}{7}{1}{2}{104}{C}
\bktocchap{2-6}{はかった\R{数}{かず}には「ずれ」がある}{\R{同}{おな}じ\R{物}{もの}を10\R{回}{かい}はかると、ぜんぶ\R{同}{おな}じになる？}{7}{1}{2}{122}{C}
}%
{「\qty{5}{m}」の「m」には、\R{計算}{けい|さん}の\R{答}{こた}えが\R{正}{ただ}しいかをたしかめる\R{力}{ちから}があります。この\R{部}{ぶ}で\R{身}{み}につける「\R{量}{りょう}＝\R{数}{かず}×\R{単位}{たん|い}」の\R{見方}{み|かた}と、マイナス・\R{分数}{ぶん|すう}・\R{割合}{わり|あい}の\R{計算}{けい|さん}は、\R{最後}{さい|ご}の\R{巻}{かん}まで\R{使}{つか}いつづけます。}
\bkpartbegin{1}{\R{量}{りょう}と\R{数}{かず}}
\bkinputchap{ch2-01}
\bkinputchap{ch2-02}
\bkinputchap{ch2-03}
\bkinputchap{ch2-04}
\bkinputchap{ch2-05}
\bkinputchap{ch2-06}

%% ======================== 第2部 大きな数・小さな数・単位 ========================
\bkPartOpener{2}{\R{第}{だい}2\R{巻}{かん}\hspace{0.8em}\R{第}{だい}2\R{部}{ぶ}}{\R{大}{おお}きな\R{数}{かず}・\R{小}{ちい}さな\R{数}{かず}・\R{単位}{たん|い}}%
  {172/{\R{紙}{かみ}を42\R{回}{かい}\R{折}{お}ると、\R{月}{つき}まで\R{届}{とど}く？},
   160/{\R{面積}{めん|せき}が\qty{2}{m^2}の\R{正方形}{せい|ほう|けい}、1\R{辺}{ぺん}は\R{何}{なん}\,m？},
   148/{\R{時速}{じ|そく}\qty{72}{km}は、\R{秒速}{びょう|そく}\R{何}{なん}\,m？}}%
  {% 倍々にのびる棒（白い線）
   \begin{scope}[white, line width=1.1pt, line join=round]
     \foreach \i/\h in {0/1.5,1/3,2/6,3/12,4/24,5/48} {\draw (124+\i*8,112) rectangle ++(6,\h);}
     \draw[line width=1.3pt, line cap=round] (118,112) -- (172,112);
   \end{scope}}%
  {とても\R{大}{おお}きな\R{数}{かず}も\R{小}{ちい}さな\R{数}{かず}も、10を\R{何回}{なん|かい}かけたかで\R{書}{か}けるよ。2\R{回}{かい}かけると2になる\R{数}{かず}や、\R{世界}{せ|かい}じゅうで\R{使}{つか}う\R{単位}{たん|い}のきまりも\R{学}{まな}ぶよ。}
\bkPartContents{2}{%% 自動生成（gen.py）
\bktocchap{2-7}{\R{何回}{なん|かい}もかける：\R{累乗}{るい|じょう}と\R{指数}{し|すう}}{\R{紙}{かみ}を42\R{回}{かい}\R{折}{お}ると、\R{月}{つき}まで\R{届}{とど}く？}{3}{1}{2}{150}{C}
\bktocchap{2-8}{2\R{回}{かい}かけると□になる\R{数}{かず}：\R{平方根}{へい|ほう|こん}}{\R{面積}{めん|せき}が2\,m²の\R{正方形}{せい|ほう|けい}、1\R{辺}{ぺん}は\R{何}{なん}\,m？}{3}{2}{2}{168}{C}
\bktocchap{2-9}{\R{分数}{ぶん|すう}で\R{書}{か}けない\R{数}{かず}がある？}{$\sqrt{2}$は、どこまで\R{計算}{けい|さん}しても\R{終}{お}わらない}{6}{2}{2}{184}{C}
\bktocchap{2-10}{\R{単位}{たん|い}の\R{世界}{せ|かい}ルール：SI\R{単位}{たん|い}}{\R{時速}{じ|そく}72\,kmは、\R{秒速}{びょう|そく}\R{何}{なん}\,m？}{7}{1}{2}{202}{C}
}%
{\R{光}{ひかり}の\R{速}{はや}さも\R{原子}{げん|し}の\R{大}{おお}きさも、\R{同}{おな}じ\R{書}{か}き\R{方}{かた}で\R{書}{か}けるようになります。\R{単位}{たん|い}をそろえる\R{力}{ちから}は、\R{計算}{けい|さん}のまちがいを\R{見}{み}つける\R{力}{ちから}になります。}
\bkpartbegin{2}{\R{大}{おお}きな\R{数}{かず}・\R{小}{ちい}さな\R{数}{かず}・\R{単位}{たん|い}}
\bkinputchap{ch2-07}
\bkinputchap{ch2-08}
\bkinputchap{ch2-09}
\bkinputchap{ch2-10}

%% ======================== 第3部 文字と式 ========================
\bkPartOpener{3}{\R{第}{だい}2\R{巻}{かん}\hspace{0.8em}\R{第}{だい}3\R{部}{ぶ}}{\R{文字}{も|じ}と\R{式}{しき}}%
  {172/{どんな\R{速}{はや}さでも\R{使}{つか}える\R{式}{しき}を、1\R{行}{ぎょう}で\R{書}{か}ける？},
   160/{1つの\R{公式}{こう|しき}から、3つの\R{公式}{こう|しき}が\R{作}{つく}れる？},
   148/{\R{先}{さき}に\R{出}{で}た\R{弟}{おとうと}に、\R{兄}{あに}は\R{何秒}{なん|びょう}で\R{追}{お}いつく？}}%
  {% 天びん（白い線）
   \begin{scope}[white, line width=1.2pt, line cap=round, line join=round]
     \draw (137,110) -- (151,110) -- (144,124) -- cycle;
     \draw (144,124) -- (144,134);
     \draw[line width=1.6pt] (121,134) -- (167,134);
     \draw[line width=0.8pt] (124,134) -- (118,121) (124,134) -- (130,121);
     \draw[line width=0.8pt] (164,134) -- (158,121) (164,134) -- (170,121);
     \draw (116,121) -- (132,121) (156,121) -- (172,121);
     \draw[rounded corners=0.8mm, line width=1pt, fill=bkVol] (119.5,121.8) rectangle (128.5,130.2);% 第7回：ひもが箱の上に見えないように、箱を色の面でぬる
     \node[text=white, font=\bfseries\large] at (124,126) {$x$};
     \fill (160.6,123.6) circle (1.8); \fill (167.4,123.6) circle (1.8); \fill (164,127.2) circle (1.8);
   \end{scope}}%
  {\R{数}{かず}のかわりに\R{文字}{も|じ}を\R{使}{つか}うと、どんな\R{数}{かず}でも\R{成}{な}り\R{立}{た}つきまりを\mbox{1\R{行}{ぎょう}で}\R{書}{か}けるよ。\R{式}{しき}を\R{天}{てん}びんのように\R{変形}{へん|けい}して、ほしい\R{量}{りょう}をとり\R{出}{だ}す\R{方法}{ほう|ほう}も\R{学}{まな}ぶよ。}
\bkPartContents{3}{%% 自動生成（gen.py）
\bktocchap{2-11}{なぜ\R{文字}{も|じ}を\R{使}{つか}うの？}{どんな\R{速}{はや}さでも\R{使}{つか}える\R{式}{しき}を、1\R{行}{ぎょう}で}{5}{2}{2}{230}{C}
\bktocchap{2-12}{\R{文字式}{も|じ|しき}のやくそく}{3\,mと2\R{秒}{びょう}は、たし\R{算}{ざん}できない}{5}{2}{2}{256}{C}
\bktocchap{2-13}{\R{公式}{こう|しき}に\R{数}{かず}を\R{入}{い}れてみよう：\R{代入}{だい|にゅう}}{\R{式}{しき}に\R{数}{かず}を\R{入}{い}れると、\R{答}{こた}えが\R{出}{で}てくる}{5}{2}{2}{272}{C}
\bktocchap{2-14}{\R{式}{しき}を\R{変形}{へん|けい}して、ほしい\R{量}{りょう}をとり\R{出}{だ}そう}{1つの\R{公式}{こう|しき}から、3つの\R{公式}{こう|しき}が\R{作}{つく}れる}{5}{2}{2}{294}{C}
\bktocchap{2-15}{\R{比例}{ひ|れい}と\R{反比例}{はん|ぴ|れい}}{ばねを2\R{倍}{ばい}のばすには、2\R{倍}{ばい}の\R{力}{ちから}}{1}{1}{2}{314}{C}
\bktocchap{2-16}{\R{何秒後}{なん|びょう|ご}に\R{追}{お}いつく？1\R{次方程式}{じ|ほう|てい|しき}}{\R{先}{さき}に\R{出}{で}た\R{弟}{おとうと}に、\R{兄}{あに}は\R{何秒}{なん|びょう}で\R{追}{お}いつく？}{1}{2}{2}{334}{C}
}%
{$x=vt$ のような\R{公式}{こう|しき}を\R{読}{よ}み、\R{数}{かず}を\R{入}{い}れ、ほしい\R{量}{りょう}について\R{解}{と}けるようになります。この\R{部}{ぶ}の「\R{式}{しき}を\R{変形}{へん|けい}する\R{力}{ちから}」は、このあとのすべての\R{巻}{かん}で\R{毎}{まい}ページ\R{使}{つか}います。}
\bkpartbegin{3}{\R{文字}{も|じ}と\R{式}{しき}}
\bkinputchap{ch2-11}
\bkinputchap{ch2-12}
\bkinputchap{ch2-13}
\bkinputchap{ch2-14}
\bkinputchap{ch2-15}
\bkinputchap{ch2-16}

%% ======================== 第4部 コンピュータで計算 ========================
\bkPartOpener{4}{\R{第}{だい}2\R{巻}{かん}\hspace{0.8em}\R{第}{だい}4\R{部}{ぶ}}{コンピュータで\R{計算}{けい|さん}}%
  {172/{0と1だけで、どんな\R{数}{かず}も\R{書}{か}ける？},
   160/{1000\R{回}{かい}の\R{計算}{けい|さん}も、コピー1\R{回}{かい}でできる？},
   148/{0.1を10\R{回}{かい}たしても、1にならない？}}%
  {% 表計算のマスと0・1（白い線）
   \begin{scope}[white, line width=1pt]
     \foreach \i in {0,...,3} \foreach \j in {0,...,3} {\draw (120+\i*12,112+\j*8) rectangle ++(12,8);}
     \foreach \i/\d in {0/1,1/0,2/1,3/1} {\node[text=white, font=\ttfamily\bfseries\large] at (126+\i*12,140) {\d};}
   \end{scope}}%
  {コンピュータは、0と1だけで\R{数}{かず}を\R{覚}{おぼ}えていて、くり\R{返}{かえ}しの\R{計算}{けい|さん}がとても\R{得意}{とく|い}だよ。\R{表計算}{ひょう|けい|さん}とプログラムで、たくさんの\R{計算}{けい|さん}をコンピュータにまかせる\R{方法}{ほう|ほう}を\R{学}{まな}ぶよ。}
\bkPartContents{4}{%% 自動生成（gen.py）
\bktocchap{2-17}{コンピュータの\R{数}{かず}：2\R{進数}{しん|すう}と16\R{進数}{しん|すう}}{0と1だけで、どんな\R{数}{かず}も\R{書}{か}ける}{3}{2}{2}{362}{C}
\bktocchap{2-18}{\R{表計算}{ひょう|けい|さん}で\R{計算}{けい|さん}しよう}{1000\R{回}{かい}の\R{計算}{けい|さん}も、コピー1\R{回}{かい}で}{5}{2}{2}{386}{C}
\bktocchap{2-19}{プログラムで\R{計算}{けい|さん}しよう}{0.1を10\R{回}{かい}たしても、1にならない？}{1}{2}{2}{406}{C}
}%
{\R{表計算}{ひょう|けい|さん}とプログラムで、くり\R{返}{かえ}しの\R{計算}{けい|さん}や、\R{表}{ひょう}からグラフを\R{作}{つく}る\R{作業}{さ|ぎょう}ができるようになります。このあとの\R{巻}{かん}では、むずかしい\R{計算}{けい|さん}をコンピュータで\R{確}{たし}かめながら\R{進}{すす}みます。}
\bkpartbegin{4}{コンピュータで\R{計算}{けい|さん}}
\bkinputchap{ch2-17}
\bkinputchap{ch2-18}
\bkinputchap{ch2-19}

%% ======================== 巻末 ========================
\bkinputchap{ch2-back}
